% !TEX root = ../algebra_02.tex

\section{Изменение координат вектора под действием линейного отображения}

\begin{Definition}{Матрица линейного отображения}{}
	Пусть $\dim V = n$, $\dim W = m$, $\alpha \in \mathrm{Hom}(V, W)$, $E = (e_1, \ldots, e_n)$~--- базис $V$, $F = (f_1, \ldots, f_m)$~--- базис $W$.
	\textbf{Матрицей отображения} $\alpha$ в базисах $E$ и $F$ называется матрица $[\alpha]_{E,F}$, столбцами которой являются координаты образов базисных векторов $e_j$ в базисе $F$:
	\[
	[\alpha]_{E,F} = \bigl([\alpha e_1]_F \;\Big|\; \ldots \;\Big|\; [\alpha e_n]_F\bigr).
	\]
	Символически это записывается как $\alpha E = F \cdot [\alpha]_{E,F}$.
\end{Definition}

Зная матрицу отображения, можно легко найти координаты образа любого вектора.

\begin{Statement}{Координаты образа вектора}{}
	Пусть $\alpha\colon V\to W$ --- линейное отображение и $A=[\alpha]_{E,F}$ --- его матрица. Тогда для любого $v\in V$:
	\[
	[\alpha v]_F = A[v]_E.
	\]
\end{Statement}

\begin{proof}
	Так как $v = E[v]_E$.
	Тогда, в силу линейности отображения $\alpha$, имеем:
	\[
	\alpha v = \alpha(E[v]_E) = (\alpha E)[v]_E.
	\]
	Подставляем $\alpha E = F A$:
	\[
	\alpha v = (FA)[v]_E = F(A[v]_E).
	\]
	 Следовательно, координатный столбец вектора $\alpha v$ равен $A[v]_E$, что и требовалось доказать.
\end{proof}
